\section{Answer to Q1: the entries of $A^k$ count walks}\label{sec:q1}

\plainwords{} Take the island table, and use as many copies of it as the trip has hops. Each cell then counts the trips with exactly that many hops from its row island to its column island. These trips may land on an island twice, so they are not all fresh-island trips.

\begin{theorem}[Walk counting]\label{thm:counting}
For every integer $k\ge 0$ and all $i,j$, the entry $(A^k)_{ij}$ is the number
of walks of length $k$ from $n_i$ to $n_j$.
\end{theorem}
The University of Galway notes state this for $k\ge 1$ with an induction
proof \cite[Theorem~1.1.3, p.~4]{galwaynotes}; the case $k=0$ is this paper's
convention. The proof is derived here.
\begin{proof}
By induction on $k$.

\emph{Base case $k=0$.} $A^0=I$. A walk of length $0$ from $n_i$ to $n_j$ is a
single vertex that is both start and end, so there is exactly one such walk
if $i=j$ and none if $i\ne j$. This is $I_{ij}$.

\emph{Inductive case.} Assume that for some $k\ge 0$ and all $i,r$ the entry
$(A^k)_{ir}$ is the number of walks of length $k$ from $n_i$ to $n_r$. Fix $i$
and $j$. Since $A^{k+1}=A^kA$,
\begin{equation}\label{eq:step}
(A^{k+1})_{ij}=\sum_{r=1}^{N}(A^k)_{ir}A_{rj}.
\end{equation}
A walk $u_0\cdots u_ku_{k+1}$ of length $k+1$ from $n_i$ to $n_j$ has
$u_k=n_r$ for exactly one $r$; sort the walks by this $r$. Fix $r$.
If $n_r$ is adjacent to $n_j$ (so $A_{rj}=1$), the walks of this kind
correspond one to one to the walks of length $k$ from $n_i$ to $n_r$: deleting
the final vertex $n_j$ gives such a walk, and appending $n_j$ to such a walk
gives a walk of the kind, because $n_r$ is adjacent to $n_j$. By the
hypothesis there are $(A^k)_{ir}=(A^k)_{ir}A_{rj}$ of them. If $n_r$ is not
adjacent to $n_j$ (so $A_{rj}=0$), no walk can end with the step from $n_r$ to
$n_j$, so there are $0=(A^k)_{ir}A_{rj}$ of them. Every walk is counted for
exactly one $r$, so the number of walks of length $k+1$ from $n_i$ to $n_j$ is
$\sum_r(A^k)_{ir}A_{rj}$, which is $(A^{k+1})_{ij}$ by \cref{eq:step}.
\end{proof}

\paragraph{Warning: walks, not paths.} The proof uses only that consecutive
vertices are adjacent, so repeated vertices are counted. For a path, appending
$n_j$ may repeat a vertex and the one to one correspondence fails. In the
board graph $(A^2)_{11}=1$ and the walk counted is $n_1n_2n_1$, which repeats
$n_1$; the only path from $n_1$ to $n_1$ is the single vertex $n_1$, of length
$0$ (a path of length $k\ge 1$ from a vertex to itself would repeat it). So
``$(A^2)_{11}$ counts the paths of length $2$ from $n_1$ to $n_1$'' is false,
although ``counts the walks'' is true. Also $(A^4)_{22}=4$ by script, so the entry is
not ``$1$ if $n_j$ can be reached from $n_i$''.

\section{Answer to Q3: the path between $n_i$ and $n_j$}\label{sec:q3}

\plainwords{} Choose two islands. A fresh-island trip between them starts on the first island and ends on the second. A frog that stays put makes such a trip from an island to itself.

\begin{definition}[Path between $n_i$ and $n_j$]\label{def:pathbetween}
A \emph{path between $n_i$ and $n_j$} (also a path \emph{from $n_i$ to $n_j$})
is a sequence $u_0,u_1,\dots,u_k$ with $k\ge 0$, $u_0=n_i$, $u_k=n_j$, in
which $u_{t-1}$ and $u_t$ are adjacent for $1\le t\le k$ and the vertices
$u_0,\dots,u_k$ are pairwise distinct. Its length is $k$.
\end{definition}
This is \cref{def:path} with the ends fixed, the textbook's
Definition~1.2.13 with ends $n_i$ and $n_j$
\cite[p.~22]{petrosyan2015graphs}. If $i\ne j$ such a path has at least one
edge; if $i=j$ the only path is $(n_i)$, since $u_0=u_k$ with $k\ge 1$ would
repeat a vertex. In the board graph $n_1n_2n_3$ is the only path from $n_1$ to $n_3$.

\section{Answer to Q2: a connectivity test from $A$}\label{sec:q2}

The test rests on two facts about walks and paths. Both are stated in
Grinberg \cite[Corollary~2.9.7, p.~39; Section~2.9.4, p.~40]{grinberg2025graph}
and the first is the textbook's Lemma~1.2.1 for $u\ne v$ without a length
bound \cite[p.~23]{petrosyan2015graphs}; the proofs are derived here.

\plainwords{} Suppose a frog gets from one island to another and lands on some island twice. Cut out all the hops between the two visits to that island. What is left is still a trip between the same two islands, with fewer hops. Cutting again and again ends in a fresh-island trip.

\begin{lemma}[Every walk contains a path]\label{lem:cut}
Let $W=u_0\cdots u_k$ be a walk from $u$ to $v$. Then there is a path from $u$
to $v$ of length at most $k$.
\end{lemma}
\begin{proof}
By strong induction on $k$: for walks of length $k$ we assume the statement
for all shorter walks. \emph{Base case $k=0$.} $W=(u)$ has one vertex, so it
is a path of length $0$. \emph{Inductive case $k\ge 1$.} If the vertices of $W$
are pairwise distinct, $W$ is a path of length $k$. Otherwise there are
positions $s<t$ with $u_s=u_t$. Let $W'=u_0\cdots u_su_{t+1}\cdots u_k$
(if $t=k$, $W'=u_0\cdots u_s$). It is a walk, because the only new
consecutive pair is $u_s,u_{t+1}$ and $u_s=u_t$ is adjacent to $u_{t+1}$. It
runs from $u$ to $v$ and has length $k-(t-s)<k$. By the hypothesis $W'$, and so
$W$, yields a path from $u$ to $v$ of length at most $k-(t-s)\le k$.
\end{proof}

\plainwords{} Pick two different islands and look at their cell in the tables for one hop, two hops, and so on. Go up to one fewer hops than there are islands. The frog can get from one island to the other exactly when one of these cells is not zero.

\begin{lemma}[Reachability from the powers]\label{lem:reach}
Let $i\ne j$. There is a path from $n_i$ to $n_j$ if and only if
$(A^k)_{ij}>0$ for some $k$ with $1\le k\le N-1$.
\end{lemma}
\begin{proof}
($\Rightarrow$) Let a path from $n_i$ to $n_j$ have length $k$. As $i\ne j$,
$k\ge 1$. A path of length $k$ has $k+1$ pairwise distinct vertices among
the $N$ vertices of the graph, so $k\le N-1$. A path is a walk, so by
\cref{thm:counting} $(A^k)_{ij}\ge 1$. ($\Leftarrow$) If $(A^k)_{ij}>0$ with
$1\le k\le N-1$, then by \cref{thm:counting} a walk of length $k$ joins $n_i$
to $n_j$, and by \cref{lem:cut} a path joins them.
\end{proof}

\plainwords{} Add the tables for one hop, two hops, and so on, up to one fewer hops than there are islands. Look only at cells for two different islands. The pond is in one piece exactly when none of these cells is zero. How large a cell is does not matter.

\begin{theorem}[Sum test]\label{thm:sum}
Let $N\ge 2$ (for $N=1$ the graph is connected) and $B=A+A^2+\dots+A^{N-1}$. Then $G$ is connected if and only if
$B_{ij}>0$ for all $i\ne j$.
\end{theorem}
This is Theorem~2 of Paramadevan and Sotheeswaran
\cite[p.~17]{paramadevan2021adjacency}; the proof here is derived.
\begin{proof}
Each $A^k$ has non-negative entries, so $B_{ij}=\sum_{k=1}^{N-1}(A^k)_{ij}$ is a
sum of non-negative terms, and it is positive exactly when at least one term
is positive. For $i\ne j$ this holds exactly when a path joins $n_i$ and
$n_j$, by \cref{lem:reach}. So $B_{ij}>0$ for all $i\ne j$ exactly when every
pair of different vertices is joined by a path, which is
\cref{def:connected}.
\end{proof}
The cutoff $N-1$ is the largest length a path can have
(proof of \cref{lem:reach}). The entries of $B$ are walk counts, not path counts; only
whether they are zero is used.
